\documentclass[11pt,a4paper]{article}

\usepackage{fontspec}
\usepackage[french]{babel}
\usepackage{microtype}
\usepackage[margin=2.25cm]{geometry}
\usepackage{amsmath,amssymb}
\usepackage{booktabs,longtable,tabularx}
\usepackage{enumitem}
\usepackage{xcolor}
\usepackage{hyperref}
\usepackage{url}
\usepackage{fancyhdr}
\usepackage{titlesec}

\setmainfont{Latin Modern Roman}
\setsansfont{Latin Modern Sans}
\setmonofont{Latin Modern Mono}

\definecolor{ink}{HTML}{18212B}
\definecolor{navy}{HTML}{173B57}
\definecolor{blue}{HTML}{176B91}
\definecolor{lightblue}{HTML}{EAF3F7}
\definecolor{rulegray}{HTML}{CFD8DE}
\definecolor{softgray}{HTML}{F5F7F8}
\definecolor{warning}{HTML}{8A4B08}

\color{ink}
\hypersetup{
  colorlinks=true,
  linkcolor=navy,
  urlcolor=blue,
  pdftitle={Prompt Fable 5 — Campagne de sensibilité M4 et M4B},
  pdfauthor={Anatole Joseph Stouls}
}
\urlstyle{tt}
\setlist[itemize]{leftmargin=1.5em,itemsep=0.3em,topsep=0.4em}
\setlist[enumerate]{leftmargin=1.7em,itemsep=0.35em,topsep=0.4em}
\setlength{\parindent}{0pt}
\setlength{\parskip}{0.55em}
\setlength{\emergencystretch}{3em}

\titleformat{\section}
  {\Large\bfseries\sffamily\color{navy}}{\thesection}{0.65em}{}
\titleformat{\subsection}
  {\large\bfseries\sffamily\color{blue}}{\thesubsection}{0.55em}{}
\titlespacing*{\section}{0pt}{1.7em}{0.55em}
\titlespacing*{\subsection}{0pt}{1.2em}{0.4em}

\pagestyle{fancy}
\setlength{\headheight}{22pt}
\fancyhf{}
\fancyhead[L]{\small\sffamily Prompt Fable 5}
\fancyhead[R]{\small\sffamily Sensibilité M4--M4B}
\fancyfoot[C]{\small\sffamily\thepage}
\renewcommand{\headrulewidth}{0.4pt}
\renewcommand{\headrule}{\hbox to\headwidth{\color{rulegray}\leaders\hrule height \headrulewidth\hfill}}

\newcommand{\code}[1]{\texttt{\color{navy}#1}}
\newcommand{\file}[1]{\nolinkurl{#1}}
\newcommand{\term}[1]{\textbf{\color{navy}#1}}
\newenvironment{briefbox}
  {\begin{center}\begin{minipage}{0.94\textwidth}\color{navy}\hrule\vspace{0.7em}}
  {\vspace{0.5em}\hrule\end{minipage}\end{center}}
\newenvironment{notebox}
  {\begin{center}\begin{minipage}{0.90\textwidth}\color{navy}\itshape}
  {\end{minipage}\end{center}}

\begin{document}

\begin{titlepage}
  \centering
  \vspace*{2.2cm}
  {\sffamily\bfseries\color{navy}\fontsize{28}{34}\selectfont
    Campagne de sensibilité\par}
  \vspace{0.35cm}
  {\sffamily\bfseries\color{blue}\fontsize{23}{29}\selectfont
    des modèles M4 et M4B\par}
  \vspace{1.1cm}
  {\Large Prompt de recherche pour Claude Fable 5\par}
  \vspace{1.6cm}
  \begin{minipage}{0.78\textwidth}
    \begin{briefbox}
      Mission de recherche numérique de bout en bout : audit des moteurs,
      plan d'expérience, simulations stochastiques, analyse de sensibilité,
      vérification et rapport scientifique reproductible.
    \end{briefbox}
  \end{minipage}
  \vfill
  {\large Anatole Joseph Stouls\par}
  \vspace{0.25cm}
  {\large 17 juillet 2026\par}
  \vspace{0.8cm}
  {\small\sffamily À exécuter de préférence avec l'effort \code{xhigh}.\par}
\end{titlepage}

\tableofcontents
\clearpage

\section{Rôle et objectif scientifique}

Tu es responsable d'une étude de sensibilité reproductible sur les modèles de
société de crédit M4 et M4B. Travaille comme chercheur en économophysique,
systèmes complexes et simulation stochastique : conçois le protocole, audite
les moteurs, exécute les simulations, analyse les résultats et rédige le
rapport final en français.

L'objectif n'est ni de chercher la configuration qui « donne le plus de SOC »,
ni de reconfirmer sélectivement le rapport M4. Il faut cartographier honnêtement :

\begin{enumerate}
  \item quels paramètres contrôlent le régime démographique, le réseau de
    crédit, les faillites en cascade et les distributions individuelles ;
  \item où se trouvent les non-linéarités, seuils, interactions et changements
    de régime ;
  \item quelles conclusions de M4 sont robustes dans une plage de paramètres
    et lesquelles dépendent d'un choix institutionnel ;
  \item si M4B reproduit bien M4 sur le sous-espace commun, sans traiter ces
    deux implémentations d'une même mécanique comme deux preuves indépendantes.
\end{enumerate}

\begin{notebox}
\textbf{Consigne d'autonomie.} Quand tu as assez d'information pour agir,
agis. Mène la mission jusqu'aux résultats et au rapport vérifié ; ne termine
pas sur un simple plan ou une promesse d'exécution.
\end{notebox}

\section{Répertoire de travail et sources obligatoires}

Travaille dans \file{/home/anatole/jupyter}. Utilise le Python du venv :
\file{/home/anatole/jupyter/.venv/bin/python3}.

Lis intégralement avant de figer le protocole :

\begin{itemize}
  \item guide Fable :
    \file{/home/anatole/jupyter/m4_credit_soc/Prompting Claude Fable.md} ;
  \item rapport scientifique M4 :
    \file{/home/anatole/jupyter/m4_credit_soc_fable/reports/01_soc_final/main.pdf}
    et sa source \code{main.tex} ;
  \item journal, configuration, moteur, statistiques et campagnes M4 :
    \file{m4_credit_soc_fable/NOTES.md}, \file{src/m4/},
    \file{experiments/m4/} et \file{tests/} ;
  \item spécification M4B : \file{m4b_credit_soc_mini/README.md} et
    \file{m4b_credit_soc_mini/report/mecanique_m4b.pdf} ;
  \item moteur et sorties M4B : \file{m4b_credit_soc_mini/m4b/},
    \file{run.py} et \file{tests/} ;
  \item intégration actuelle de M4B dans Simulation Lab :
    \file{modeles-systeme-physicoeconomique/m4b_credit_soc_mini/} et
    \file{simulation_lab/}.
\end{itemize}

Crée tous les scripts, résultats, figures, notes et rapports de l'étude dans :

\begin{center}
  \file{/home/anatole/jupyter/recherche/sensibilite_m4_m4b/}
\end{center}

Traite comme \term{sources en lecture seule} les moteurs
\file{m4_credit_soc_fable/src/} et \file{m4b_credit_soc_mini/m4b/}, leurs
tests, leurs rapports existants et tous les résultats antérieurs. Ne modifie
pas la dynamique pour faciliter l'étude. Une lacune d'instrumentation doit
être contournée dans les scripts de campagne ou signalée ; elle ne justifie
pas silencieusement une nouvelle version du modèle.

\section{Relation exacte entre M4 et M4B}

Ne pars pas de l'hypothèse qu'il s'agit de deux modèles concurrents. M4B est
une réduction indépendante qui reproduit la configuration finale de M4. Sur
leur sous-espace commun, la trajectoire doit être identique à graine et
configuration identiques.

\begin{longtable}{@{}p{0.18\textwidth}p{0.72\textwidth}@{}}
\toprule
\textbf{Paramètre} & \textbf{Rôle scientifique dans M4B} \\
\midrule
\endhead
\code{lam} & Intensité des naissances et axe de taille du système. \\
\code{delta} & Dépréciation du capital. \\
\code{sigma} & Volatilité du choc multiplicatif. \\
\code{K0} & Capital de naissance. \\
\code{k} & Taille de l'échantillon du marché local. \\
\bottomrule
\end{longtable}

\code{seed} et \code{T} définissent la réplication et l'horizon expérimental.
\code{pop\_max} est un garde-fou d'arrêt, pas un mécanisme de bornage.
\code{alpha=1} et les tolérances numériques sont des conventions ou constantes
et ne doivent pas être promues artificiellement en paramètres scientifiques.

M4 seul expose les choix structurels retirés de M4B, notamment \code{credit},
\code{objective}, \code{rounds\_div}, les variantes de naissance, la
corrélation des chocs et les règles de faillite. Utilise M4 uniquement pour
cette étude de robustesse institutionnelle et pour quelques contrôles de
parité. Utilise M4B pour la campagne quantitative sur les paramètres communs.
\term{Ne double pas tous les runs.}

\section{Questions et grandeurs de réponse}

Définis avant les runs un dictionnaire des métriques, leur formule, leur unité,
leur fenêtre temporelle, leurs conditions de validité et leur statut primaire
ou secondaire.

\subsection{Régime et démographie}

\begin{itemize}
  \item statut \code{ok}, extinction ou arrêt par \code{pop\_max} ;
  \item population moyenne après burn-in, $N/\lambda$, pente temporelle et
    variabilité ;
  \item naissances, morts, durée de vie et renouvellement ;
  \item stationnarité vérifiée sur plusieurs fenêtres, pas seulement sur la
    valeur finale.
\end{itemize}

\subsection{Crédit, exposition et réseau}

\begin{itemize}
  \item nombre et volume de prêts, dette rapportée au capital ou aux actifs,
    degrés et concentration des expositions ;
  \item rythme d'accumulation puis d'élagage du carnet autour des grands
    événements ;
  \item métriques normalisées par la population lorsque nécessaire.
\end{itemize}

\subsection{Avalanches et propagation causale}

\begin{itemize}
  \item fréquence, quantiles, maximum absolu et maximum rapporté à la
    population ;
  \item rapport de branchement, fraction induite, racines/taille et profondeur ;
  \item susceptibilité ou échelle de coupure, par exemple
    $\langle s^2\rangle/\langle s\rangle$ ;
  \item MLE \term{discret} de la loi de puissance, notamment au seuil
    $s_{\min}=2$, ajustement loi de puissance avec coupure et comparaison de
    vraisemblance avec une log-normale discrète renormalisée ;
  \item scaling de la coupure avec la taille du système.
\end{itemize}

Ne conclus jamais à une loi de puissance sur le seul $r^2$ d'une régression
log-log. Si un run contient trop peu d'avalanches pour un ajustement fiable,
marque la métrique non identifiable au lieu de produire un nombre trompeur.

\subsection{Distributions individuelles et inégalités}

\begin{itemize}
  \item capital $K$, valeur nette, revenu brut/net, Gini et renouvellement du
    haut de la distribution ;
  \item familles de distribution comparées avec les mêmes supports et les
    mêmes règles de troncature ;
  \item contrôle par âge ou cohorte lorsque l'interprétation distributionnelle
    peut en dépendre.
\end{itemize}

\subsection{Intégrité numérique}

\begin{itemize}
  \item égalité créances-dettes, absence de contrats orphelins, capital non
    négatif, bilan réel par pas et reproductibilité ;
  \item temps de calcul, volume des données et tout arrêt anticipé.
\end{itemize}

\section{Protocole expérimental}

\subsection{Audit, parité et réutilisation}

\begin{enumerate}
  \item Exécute les suites de tests M4 et M4B.
  \item Vérifie au moins deux trajectoires M4/M4B sur des configurations
    communes non triviales. Compare les séries par pas, états, contrats et
    avalanches, pas seulement les moyennes finales.
  \item Inventorie les runs existants. Ne réutilise un run que si son moteur,
    sa configuration complète, sa graine, son horizon, son statut et ses
    fichiers primaires sont vérifiables. Un nom de dossier ne suffit pas.
  \item Lance un petit benchmark représentatif pour mesurer coût et taille
    disque, puis fixe un budget expérimental défendable. Laisse au moins deux
    cœurs au système et évite la surallocation mémoire ou disque.
\end{enumerate}

\subsection{Séparer les axes qui n'ont pas le même sens}

Traite \code{lam} d'abord comme axe de taille finie et démographique, et
\code{T} comme axe de convergence temporelle. Ne les mélange pas sans
précaution dans un unique classement d'importance avec \code{delta},
\code{sigma}, \code{K0} et \code{k}.

La référence est :

\begin{center}
\code{lam=10, delta=0.05, sigma=0.25, K0=25, k=3, T=2000}.
\end{center}

Pour obtenir des queues mieux identifiées, le centre de la campagne peut être
\code{lam=30} ; justifie ce choix et conserve le lien avec la baseline.

Construis successivement :

\begin{enumerate}
  \item un contrôle de taille et d'horizon autour de
    $\lambda\in\{10,30,100\}$ et $T\in\{2000,4000\}$, en réutilisant les
    runs M4 valides déjà présents ;
  \item des courbes OAT autour du centre pour rendre visibles monotonies,
    seuils et formes non linéaires ;
  \item un plan global espace-remplissant sur les paramètres microscopiques
    communs ;
  \item des coupes 2D ciblées seulement après le screening ;
  \item une phase confirmatoire indépendante sur les régimes et interactions
    les plus importants.
\end{enumerate}

\subsection{Plages initiales}

Ces plages doivent être auditées par des pilotes, puis corrigées si elles ne
produisent que des cas triviaux ou des arrêts :

\begin{center}
\begin{tabular}{@{}lll@{}}
\toprule
\textbf{Paramètre} & \textbf{Plage initiale} & \textbf{Remarque} \\
\midrule
\code{delta} & $0{,}02$--$0{,}10$ & autour de $0{,}05$ \\
\code{sigma} & $0{,}10$--$0{,}50$ & contrôle $\sigma=0$ séparé \\
\code{K0} & $5$--$100$ & échelle logarithmique préférable \\
\code{k} & $\{2,3,4,6,10\}$ & paramètre discret \\
\code{lam} & $\{10,30,100\}$ & scaling principal \\
\bottomrule
\end{tabular}
\end{center}

Ces bornes sont des points de départ, pas des vérités. Relie-les aux équations,
au rapport, aux pilotes et aux régimes observés. Ne fais pas varier
\code{pop\_max} comme s'il bornait la population ; augmente-le seulement pour
diagnostiquer un arrêt, sans masquer un régime explosif.

Pour le plan global, utilise un plan reproductible adapté au mélange continu
et discret, par exemple un Latin hypercube stratifié pour les continus avec
équilibrage de \code{k}. Un ordre de grandeur de 48 à 80 points avec trois
graines par point est raisonnable après benchmark ; réduis ou augmente ce
nombre selon le coût et la précision observés, en documentant la décision. Ne
revendique pas des indices de Sobol si le plan de Saltelli et le nombre de
réplications ne les rendent pas valides.

Utilise le même panel de graines pour chaque cellule d'une comparaison. Cela
permet des contrastes appariés, mais ne l'appelle pas abusivement « common
random numbers » si les trajectoires consomment ensuite des tirages différents.
Ne fusionne jamais les événements de plusieurs graines avant d'ajuster une
loi : ajuste chaque graine, puis agrège les estimateurs et leur incertitude.

Après screening, confirme les cellules retenues avec $T=4000$ et au moins cinq
graines. Près d'une frontière de régime, ou lorsque la variance inter-graines
domine l'effet, augmente les répétitions plutôt que d'affirmer une sensibilité
instable. Utilise par défaut un burn-in de $T/4$ et vérifie la robustesse à la
fenêtre.

\subsection{Robustesse institutionnelle dans M4}

Sépare clairement cette partie de la sensibilité paramétrique M4B. Étudie de
façon parcimonieuse et causale :

\begin{itemize}
  \item la dose-réponse de \code{rounds\_div} dans les valeurs supportées par
    le moteur, notamment $\{6,3,2,1\}$ ;
  \item \code{objective="income"} contre \code{"wealth"} ;
  \item la combinaison factorielle des règles de faillite
    \code{cancel/transfer} $\times$ \code{destroy/prorata} ;
  \item crédit actif/inactif, naissance par emprunt et chocs corrélés comme
    contrôles structurels, si les rapports ou runs existants ne suffisent pas ;
  \item \code{k} aux mêmes points dans M4 et M4B sur un petit sous-échantillon
    de parité.
\end{itemize}

Réutilise en priorité les ablations M4 existantes dont la provenance est
complète. Distingue les résultats obtenus autour de \code{rounds\_div=3} et de
la configuration finale \code{rounds\_div=1} ; ne transpose pas
quantitativement une ablation d'un centre à l'autre sans contrôle.

\subsection{Analyse de sensibilité}

Produis plusieurs lectures complémentaires :

\begin{itemize}
  \item courbes OAT avec incertitude inter-graines et contrastes appariés ;
  \item corrélations de rang partielles ou régressions standardisées avec
    intervalles bootstrap pour le screening global ;
  \item modèle de réponse non linéaire validé hors échantillon pour détecter
    seuils et interactions, avec importance par permutation si pertinente ;
  \item cartes de régime extinction/stationnaire/explosion et diagrammes 2D
    des interactions confirmées ;
  \item décomposition explicite de la variabilité due aux paramètres et de la
    variabilité stochastique inter-graines.
\end{itemize}

Ne donne pas un classement unique sans préciser la métrique : un paramètre
peut être peu important pour le rapport de branchement mais dominant pour la
population ou les inégalités. Rapporte tailles d'effet, incertitudes,
non-monotonies et conditions de validité. Toute analyse exploratoire utilisée
pour choisir des cellules doit être étiquetée comme telle ; les conclusions
fortes doivent venir de la phase confirmatoire.

\section{Infrastructure, traçabilité et Simulation Lab}

Organise le dossier de campagne au minimum ainsi :

\begin{verbatim}
recherche/sensibilite_m4_m4b/
  README.md
  design.md
  notes.md
  scripts/
  manifests/
  results/
  figures/
  report/
\end{verbatim}

Exigences :

\begin{itemize}
  \item scripts relançables et reprise après interruption ;
  \item identifiant de run déterministe ou table de correspondance non ambiguë ;
  \item manifeste comportant modèle, version ou provenance, tous les
    paramètres, graine, horizon, burn-in, statut, durée, chemin des données et
    contrôles d'intégrité ;
  \item résultats agrégés en CSV/JSON lisibles sans relancer les simulations ;
  \item journal des décisions, résultats négatifs et anomalies dans
    \code{notes.md} ;
  \item aucune cellule silencieusement exclue parce qu'elle s'éteint, explose
    ou contredit l'hypothèse.
\end{itemize}

M4B est le modèle actif de Simulation Lab. Les runs M4B confirmatoires retenus
doivent y être consultables, avec leurs figures et métadonnées, et être marqués
à conserver. Pour le screening massif, limite l'écriture individuelle
(\code{individual\_every=0} ou fréquence espacée) et les instantanés lorsque
ces données ne sont pas requises ; cela ne doit pas changer la trajectoire.
Pour les cellules confirmatoires et distributionnelles, conserve les mesures
nécessaires. Vérifie toujours cette neutralité de mesure sur un contrôle.

N'essaie pas de réactiver M4 dans Simulation Lab : les ablations M4 peuvent
être exécutées par les outils propres à \file{m4_credit_soc_fable} et
référencées dans le manifeste commun.

\section{Vérification et conduite autonome}

Délègue les sous-tâches réellement indépendantes à des sous-agents, par exemple
l'audit de parité, l'inventaire des résultats existants et la vérification
statistique. Garde la conception expérimentale et la synthèse scientifique sous
ta responsabilité. À la fin de chaque phase majeure, fais vérifier par un
sous-agent à contexte frais : conformité au protocole, fuites entre exploration
et confirmation, calcul des métriques, provenance des figures et adéquation des
conclusions aux données.

Avant chaque compte rendu de progression, vérifie chaque affirmation dans un
résultat d'outil de cette session. Si un test échoue, si un run est incomplet ou
si une étape est omise, dis-le explicitement. Ne présente jamais une intention
comme un travail accompli.

Travaille de façon autonome pour toutes les actions réversibles dans ce
périmètre. Ne demande l'utilisateur que si une action destructive ou
irréversible, un vrai changement de modèle, une extension de périmètre ou une
information que lui seul possède devient indispensable. Ne supprime aucun
résultat existant et ne modifie pas les moteurs de référence.

\section{Livrables finaux}

\begin{enumerate}
  \item \code{design.md} : protocole figé avant la campagne confirmatoire,
    domaines, métriques, graines, critères d'arrêt et plan d'analyse.
  \item Manifeste complet des runs et table longue des métriques par run et
    graine.
  \item Tables d'effets OAT, résultats du screening global, interactions et
    incertitudes.
  \item Figures lisibles : courbes de réponse, diagrammes de régime,
    interactions, sensibilité par famille de métriques, scaling de taille
    finie et contrôles de parité.
  \item Rapport scientifique final en français, en \code{.tex} et
    \code{.pdf}, contenant méthode, résultats, résultats négatifs, comparaison
    M4/M4B, limites et conclusions robustes. Chaque figure doit être traçable
    jusqu'aux runs.
  \item \code{README.md} avec les commandes exactes pour reproduire l'étude et
    régénérer les figures ou le rapport sans relancer inutilement les
    simulations.
  \item Une synthèse courte répondant explicitement aux questions suivantes :
    \begin{itemize}
      \item quels paramètres influencent quelles sorties ;
      \item où sont les frontières de régime et interactions ;
      \item quelles signatures SOC sont robustes ;
      \item ce qui relève d'un effet de taille, d'un choix institutionnel ou
        du bruit stochastique ;
      \item ce que la parité M4/M4B permet, et ne permet pas, de conclure.
    \end{itemize}
\end{enumerate}

\begin{briefbox}
\textbf{Définition de « terminé ».} La mission est terminée seulement lorsque
les runs retenus sont exécutés ou explicitement comptabilisés comme échecs,
les contrôles sont passés, les résultats sont reproductibles, les figures sont
sourcées et le PDF final est compilé. Dans ton message final, commence par le
résultat scientifique le plus important, puis donne les chemins des livrables
et les éventuelles limites.
\end{briefbox}

\end{document}
